\section{Problem Definition}
\label{sec:problem-definition}

Nearly every machine learning model, graph neural networks (GNNs) included, exposes a set of hyperparameters that must be fixed prior to training. Selecting good values for them is the task of hyperparameter optimization (HPO). Although such tuning was long carried out by hand, the field has come to accept that principled, automated HPO yields models that both perform and generalize better. From an optimization standpoint, HPO is often a black-box problem, and an unusually costly one: assessing a single configuration entails training a model and measuring its validation performance.

Our focus throughout is HPO for supervised learning, and we adopt the notation and formalism introduced by \cite{bischl_hyperparameter_2023}. Let \( \mathcal{D} \in \mathrm{D} \) be a labelled dataset made up of pairs \( \left( \mathvec{x}_i, y_i \right) \) with features \( \mathvec{x}_i \in \mathcal{X} \) and labels \( y_i \in \mathcal{Y} \), and take its elements to be i.i.d.\ draws from an unknown joint distribution \( \mathrm{P}_{\mathcal{X}\mathcal{Y}} \). When it is convenient to treat the dataset as a whole, we stack its features into the matrix \( \mathmat{X}_{\mathcal{D}} \) and its labels into the vector \( \mathvec{y}_{\mathcal{D}} \).

Supervised learning seeks a function \( f: \mathcal{X} \to \mathfield{R}^g \) that turns features into predictions and behaves well on data it has not seen, where \( g \) equals the number of classes in classification and \( g = 1 \) in regression. The collection of admissible such functions is the \name{hypothesis space} \( \mathscr{H} \). Producing a concrete model depends on a hyperparameter configuration (i.\ e.\ a vector of all hyperparameters) \( \lambda \in \boldsymbol\Lambda \), where \( \boldsymbol\Lambda \) is the \name{search space}, and we capture the training procedure as an \name{inducer} \( \mathscr{F} : \mathrm{D} \times \boldsymbol\Lambda \to \mathscr{H} \). Given a dataset \( \mathcal{D} \in \mathrm{D} \) and a configuration \( \lambda \in \boldsymbol\Lambda \), the inducer returns a fitted model \( \hat{f} = \mathscr{F} ( \mathcal{D}, \lambda ) \), which we occasionally write as \( \hat{f}_\lambda \) to stress its dependence on \( \lambda \); here \( \mathrm{D} = \bigcup_{n \in \mathfield{N}} \left( \mathcal{X} \times \mathcal{Y} \right)^n \) denotes the space of all datasets.

Once a model \( \hat{f} \) has been fitted, its quality is judged on fresh data through a performance measure \( \rho : \bigcup_{m \in \mathfield{N}} \left( \mathcal{Y}^m \times \mathfield{R}^{g \times m} \right) \to \mathfield{R} \). Given the true labels and the corresponding predictions over \( m \) samples, \( \rho \) returns a scalar score -- such as accuracy in a classification setting or mean squared error in a regression one, to name two common choices.

Since the inducer \( \mathscr{F} \) is governed by its configuration \( \lambda \in \boldsymbol\Lambda \), HPO amounts to locating the configuration \( \lambda^* \) that maximizes \( \rho \) on the held-out part of \( \mathcal{D} \). In practice, computational limits restrict the search to a finite subset \( \tilde{\Lambda} \subset \boldsymbol\Lambda \), the \name{search space}. A typical HPO algorithm proceeds iteratively: it picks configurations from \( \tilde{\Lambda} \), trains the corresponding models, scores them with \( \rho \), and uses those scores to decide which configurations to try next. We formalize such an algorithm as a tuner \( \tau : \left( \mathcal{D}, \mathscr{F}, \tilde{\Lambda}, \rho \right) \to \hat\lambda \), which, given a dataset \( \mathcal{D} \), an inducer \( \mathscr{F} \), a search space \( \tilde{\Lambda} \), and a performance measure \( \rho \), outputs an estimate \( \hat{\lambda} \in \tilde{\Lambda} \) of the optimum \( \lambda^* \).
